\ifdefined\savedparindent\else\newlength{\savedparindent}\fi
\ifdefined\savedparskip\else\newlength{\savedparskip}\fi
\setlength{\savedparindent}{\parindent}
\setlength{\savedparskip}{\parskip}
\setlength{\parindent}{0pt}
\setlength{\parskip}{0.6\baselineskip plus 2pt}

\section{Formal Scheduling Problem}
\label{sec:problem}

The following definitions use the terminology
of the project problem description~\cite{problem} and are
shared by all three search models.

\subsection{Components, slots, and domains}
\label{sec:components}

\paragraph{Components.}
Let $\mathit{Lectures}$ contain the lecture sections and
$\mathit{Tutorials}$ contain the tutorials and labs, which
are treated identically. The components to schedule are
\[
X=\mathit{Lectures}+\mathit{Tutorials},
\]
where $+$ denotes the union of these disjoint sets.
For $x\in X$, $\mathit{course}(x)$ is the tuple containing
its program and course number, such as $(\texttt{CPSC},433)$.

A tutorial with a lecture identifier belongs to that section;
a tutorial without one is shared by all sections of its course.
Define
\begin{equation}
\mathit{Lec}(t)=
\begin{cases}
\{\text{the lecture section named in }t\},
  & \text{if }t\text{ has a \texttt{LEC} identifier},\\
\{l\in\mathit{Lectures}:
  \mathit{course}(l)=\mathit{course}(t)\},
  & \text{otherwise}.
\end{cases}
\label{eq:parent}
\end{equation}

\paragraph{Weekly slots.}
A slot is a tuple
\begin{equation}
s=(\mathit{days}(s),\mathit{start}(s),\mathit{end}(s)).
\label{eq:slot}
\end{equation}
Assigning a component to $s$ schedules it on every day in
$\mathit{days}(s)$ at the stated time. This enforces the
required links between weekdays.

We use M, T, W, R, and F for Monday through Friday.
The available slot patterns are shown in
Table~\ref{tab:slots}.

\begin{table}[ht]
\centering
\small
\caption{Weekly slot patterns and permitted start times.}
\label{tab:slots}
\begin{tabular}{@{}llll@{}}
\toprule
Type & Days & Duration & Start times\\
\midrule
Lecture & MWF & 1 hour
  & 8:00, 9:00, \ldots, 20:00\\
Lecture & TR & 1.5 hours
  & 8:00, 9:30, \ldots, 18:30\\
Tutorial & MW & 1 hour
  & 8:00, 9:00, \ldots, 20:00\\
Tutorial & TR & 1 hour
  & 8:00, 9:00, \ldots, 20:00\\
Tutorial & F & 2 hours
  & 8:00, 10:00, \ldots, 18:00\\
\bottomrule
\end{tabular}
\end{table}

Let $\mathit{Slots}_L$ and $\mathit{Slots}_T$ contain the
available lecture and tutorial slots, respectively, and let
$\mathit{Slots}=\mathit{Slots}_L+\mathit{Slots}_T$.
Each slot has specified room capacities and desired minimum
occupancies. Lectures use only lecture slots, and tutorials
use only tutorial slots.

\paragraph{Time overlap.}
Two slots overlap when they share a day and their time
intervals intersect:
\begin{equation}
\mathit{overlap}(s, s') \iff
\mathit{days}(s) \cap \mathit{days}(s') \neq \varnothing
\ \wedge\ \mathit{start}(s) < \mathit{end}(s')
\ \wedge\ \mathit{start}(s') < \mathit{end}(s).
\label{eq:overlap}
\end{equation}
For example, a TR lecture from 9:30 to 11:00 overlaps a
TR tutorial from 10:00 to 11:00. Meetings that end exactly
when another begins do not overlap.

Under the specified patterns, slots of the same type
overlap exactly when they are equal:
\begin{equation}
\mathit{overlap}(s,s')\iff s=s'
\quad
\text{if }s,s'\in\mathit{Slots}_L
\text{ or }s,s'\in\mathit{Slots}_T.
\label{eq:overlap-eq}
\end{equation}

\paragraph{Domains and assignments.}
The domain $D(x)$ contains slots of the correct type for
$x$ that satisfy its individual restrictions. These include
preassignments, unwanted slots, evening requirements,
active-learning availability, and the department meeting
restriction. A preassignment restricts $D(x)$ to the
specified slot if that slot meets the other restrictions;
otherwise, $D(x)=\varnothing$.

A complete assignment is a function
\[
\mathit{assign}:X\to\mathit{Slots},
\qquad
\mathit{assign}(x)\in D(x)\quad\text{for every }x\in X.
\]
A partial assignment is defined only on the components
in $\operatorname{dom}(\mathit{assign})\subseteq X$.
This differs from $\mathit{partassign}$, which records
mandatory assignments supplied with the instance.

\paragraph{Slot counts.}
For the assignment under consideration, define
\begin{align*}
n_L(s)
 &=|\{l\in\mathit{Lectures}\cap
       \operatorname{dom}(\mathit{assign}):
       \mathit{assign}(l)=s\}|,\\
n_T(s)
 &=|\{t\in\mathit{Tutorials}\cap
       \operatorname{dom}(\mathit{assign}):
       \mathit{assign}(t)=s\}|.
\end{align*}
The counts $n_{L,\mathrm{AL}}(s)$ and
$n_{T,\mathrm{AL}}(s)$ are defined in the same way, but
include only components with an $\mathit{activelearning}$
requirement. These components also count toward the total
lecture or tutorial capacity.

\subsection{Hard constraints}
\label{sec:hard}
Hard constraints are the rules a valid timetable cannot break. Besides the sets defined above, they use these inputs:
\begin{itemize}
  \item lecture capacities $\mathit{lecturemax}(s)$ and $\mathit{ALlecturemax}(s)$ for each $s \in \mathit{Slots}_L$;
  \item tutorial capacities $\mathit{tutorialmax}(s)$ and $\mathit{ALtutorialmax}(s)$ for each $s \in \mathit{Slots}_T$;
  \item the $\mathit{notcompatible}(x,y)$ and $\mathit{unwanted}(x,s)$ statements;
  \item the preassignment $\mathit{partassign} : \mathit{Lectures} + \mathit{Tutorials} \to \mathit{Slots} + \{\$\}$, where $\$$ means no preassignment;
  \item the $\mathit{activelearning}(x)$ statements.
\end{itemize}

Here $\mathit{ALmax}(s)$ means $\mathit{ALlecturemax}(s)$ if $s \in \mathit{Slots}_L$ and $\mathit{ALtutorialmax}(s)$ if $s \in \mathit{Slots}_T$. These constraints also apply while the timetable is being built.
We check room capacities and time restrictions using the components
already assigned to slots. A conflict between two components is
checked once both have been assigned.
\begin{align}
C_1 &:\ \forall s \in \mathit{Slots}_L:\ n_L(s) \le \mathit{lecturemax}(s) && \text{(lecture capacity)}\\
C_2 &:\ \forall s \in \mathit{Slots}_T:\ n_T(s) \le \mathit{tutorialmax}(s) && \text{(tutorial capacity)}\\
C_3 &:\ \forall s \in \mathit{Slots}_L:\ n_{L,\mathit{AL}}(s) \le \mathit{ALlecturemax}(s) && \text{(AL lecture capacity)}\\
C_4 &:\ \forall s \in \mathit{Slots}_T:\ n_{T,\mathit{AL}}(s) \le \mathit{ALtutorialmax}(s) && \text{(AL tutorial capacity)}\\
C_5 &:\ \forall x \text{ with } \mathit{activelearning}(x):\ \mathit{ALmax}(\mathit{assign}(x)) > 0 && \text{(AL placement)}\\
C_6 &:\ \forall t \in \mathit{Tutorials},\ \forall l \in \mathit{Lec}(t):\ \neg\,\mathit{overlap}(\mathit{assign}(l), \mathit{assign}(t)) && \text{(lecture/tutorial)}\\
C_7 &:\ \forall\, \mathit{notcompatible}(x,y):\ \neg\,\mathit{overlap}(\mathit{assign}(x), \mathit{assign}(y)) && \text{(not compatible)}\\
C_8 &:\ \forall x \text{ with } \mathit{partassign}(x) \neq \$:\ \mathit{assign}(x) = \mathit{partassign}(x) && \text{(preassignment)}\\
C_9 &:\ \forall\, \mathit{unwanted}(x,s):\ \mathit{assign}(x) \neq s && \text{(unwanted)}\\
C_{10} &:\ \forall x \in X_9:\ \mathit{start}(\mathit{assign}(x)) \ge 18{:}00 && \text{(LEC 9 evening)}\\
C_{11} &:\ \forall l,l'\in L_5,\ l\neq l':
\ \neg\,\mathit{overlap}(\mathit{assign}(l),\mathit{assign}(l'))
&& \text{(5XX non-overlap)}\\
C_{12} &:\ \forall l\in\mathit{Lectures}:
\ \mathit{assign}(l)\neq \mathrm{TR}\ 11{:}00
&& \text{(department meeting)}\\
C_{13} &:\ \text{special CPSC 851/913 bookings (below)} && \notag
\end{align}

$X_9$ contains components whose name includes a \texttt{LEC}
number beginning with 9, including tutorials explicitly attached
to such a section. A shared tutorial without a \texttt{LEC}
identifier is not included solely because its course has an
evening section. $L_5$ contains all 500-level lecture sections
across programs. Lecture--tutorial conflicts are checked using
time overlap because lecture and tutorial slots follow
different time grids.

\textbf{Special bookings ($C_{13}$):}
If a CPSC 351 lecture is requested, we include one special
CPSC 851 tutorial/lab booking at TR 18:00--19:00 and add
incompatibilities with every CPSC 351 lecture, tutorial,
and lab. A CPSC 413 lecture similarly triggers one CPSC 913
booking at TR 18:00--19:00, incompatible with every
CPSC 413 component.

These bookings are included before search begins, so their
fixed times and incompatibilities are enforced through
$C_8$ and $C_7$. We assume that each booking uses one unit
of tutorial capacity in its assigned slot.

\textbf{Validity:}
Let $\mathcal C$ be the collection of hard-constraint
conditions described above. For a complete assignment
using the permitted domains, define
\begin{equation}
\mathrm{Constr}(\mathit{assign})
=\bigwedge_{C\in\mathcal C}C(\mathit{assign}).
\label{eq:const}
\end{equation}
Thus $\mathrm{Constr}(\mathit{assign})=\mathsf{true}$
exactly when every hard constraint holds.

A partial assignment is consistent when all assigned
components use permitted slots and their assignments
violate no hard constraint:
\begin{equation}
\begin{aligned}
\mathrm{Consistent}(\mathit{assign})\iff{}&
\bigl(\forall x\in\operatorname{dom}(\mathit{assign}):
      \mathit{assign}(x)\in D(x)\bigr)\\
&{}\land\text{no hard constraint is violated by}\\
&\phantom{{}\land{}}\text{the assignments made so far.}
\end{aligned}
\label{eq:consistent}
\end{equation}

Consistency does not guarantee that the remaining components
can be scheduled. However, a violation cannot be repaired
by adding assignments while keeping existing assignments
fixed. The tree searches can therefore discard such a branch.

\subsection{Soft constraints and objective}
\label{sec:soft}

Soft constraints measure the quality of a complete valid
timetable. They use the following inputs:
\begin{itemize}
  \item desired minimum occupancies
  $\mathit{lecturemin}(s)$ and $\mathit{tutorialmin}(s)$;
  \item preference values
  $\mathit{preference}:X\times\mathit{Slots}\to\mathbb{N}_0$,
  with zero where no preference is stated;
  \item a list $\mathcal P$ of $\mathit{pair}(x,y)$ requests;
  \item nonnegative penalty values
  $\mathit{pen}_{\mathit{lecturemin}}$,
  $\mathit{pen}_{\mathit{tutorialmin}}$,
  $\mathit{pen}_{\mathit{notpaired}}$, and
  $\mathit{pen}_{\mathit{section}}$.
\end{itemize}
Here $\mathbb{N}_0$ denotes the nonnegative integers.
The total evaluation consists of the following five
penalty contributions.

\paragraph{Lecture minimums.}
For each slot, every lecture missing from its desired minimum
adds $\mathit{pen}_{\mathrm{lecturemin}}$. A slot at or above
its minimum contributes zero:
\begin{equation}
E_{L\min}(\mathit{assign})
=\mathit{pen}_{\mathit{lecturemin}}
 \sum_{s\in\mathit{Slots}_L}
 \max(0,\mathit{lecturemin}(s)-n_L(s)).
\label{eq:elmin}
\end{equation}

\paragraph{Tutorial minimums.}
Similarly, each tutorial below a slot's desired minimum
adds $\mathit{pen}_{\mathit{tutorialmin}}$:
\begin{equation}
E_{T\min}(\mathit{assign})
=\mathit{pen}_{\mathit{tutorialmin}}
 \sum_{s\in\mathit{Slots}_T}
 \max(0,\mathit{tutorialmin}(s)-n_T(s)).
\label{eq:etmin}
\end{equation}

\paragraph{Time preferences.}
For each component, we add the preference values of all
preferred slots it does not receive:
\begin{equation}
E_{\mathrm{pref}}(\mathit{assign})
=\sum_{x\in X}
 \sum_{\substack{s\in\mathit{Slots}\\
                 s\neq\mathit{assign}(x)}}
 \mathit{preference}(x,s).
\label{eq:epref}
\end{equation}

\paragraph{Requested pairs.}
As clarified for this project, each pairing request
contains two lectures or two tutorials. Each request
whose components occupy different slots adds
$\mathit{pen}_{\mathit{notpaired}}$:
\begin{equation}
E_{\mathrm{pair}}(\mathit{assign})
=\mathit{pen}_{\mathit{notpaired}}
 \sum_{(x,y)\in\mathcal P}
 \mathbf{1}[\mathit{assign}(x)\neq\mathit{assign}(y)].
\label{eq:epair}
\end{equation}
Here $\mathbf{1}[Q]$ equals 1 when $Q$ is true and 0
otherwise. The sum counts each supplied pairing request.

\paragraph{Lecture sections sharing a slot.}
Let $\mathcal R$ contain the unordered pairs of distinct
lecture sections belonging to the same course:
\[
\mathcal R=
\{\{l,l'\}:l,l'\in\mathit{Lectures},\
l\neq l',\
\mathit{course}(l)=\mathit{course}(l')\}.
\]
Each such pair assigned to the same slot adds
$\mathit{pen}_{\mathit{section}}$:
\begin{equation}
E_{\mathrm{section}}(\mathit{assign})
=\mathit{pen}_{\mathit{section}}
 \sum_{\{l,l'\}\in\mathcal R}
 \mathbf{1}[\mathit{assign}(l)=\mathit{assign}(l')].
\label{eq:esection}
\end{equation}
Using unordered pairs counts each pair of sections once.

\paragraph{Objective.}
The evaluation function is the sum of these five penalties:
\begin{equation}
\begin{aligned}
\mathrm{Eval}(\mathit{assign})={}&
 E_{L\min}(\mathit{assign})
 +E_{T\min}(\mathit{assign})\\
&+E_{\mathrm{pref}}(\mathit{assign})
 +E_{\mathrm{pair}}(\mathit{assign})\\
&+E_{\mathrm{section}}(\mathit{assign}).
\end{aligned}
\label{eq:eval}
\end{equation}
We seek a complete assignment that minimizes this total
while satisfying every hard constraint:
\[
\min_{\mathit{assign}}\mathrm{Eval}(\mathit{assign})
\quad\text{subject to}\quad
\mathrm{Constr}(\mathit{assign})=\mathsf{true}.
\]
All penalties are nonnegative, so a valid timetable with
evaluation zero is globally optimal.